\uuid{3UH2}
\titre{Logique}
\auteur{bodin ; croizat ; sacre}
\organisation{exo7}
\licence{CC BY-NC-SA 4.0}
\datecreate{2022-03-08}
\niveau{L1}
\module{Algèbre}
\chapitre{Logique, ensemble, raisonnement}
\sousChapitre{Logique}
\difficulte{2}
\theme{Logique}
\exo7qcmid{hilisit.logique.008}

\contenu{
\begin{question}
\qclassification{100.01}
Soit $x$ un réel. On considère l'implication :
$$\text{"}x^2>0\implies x>0\text{".}$$
Quelles sont les affirmations vraies ?
\begin{answers}
	\good{L'implication réciproque est "$x>0\implies x^2>0$".}
	\bad{La contraposée est "$x>0\implies x^2>0$".}
	\bad{Si l'implication est fausse alors l'implication réciproque l'est aussi.}
	\good{Si l'implication est fausse alors sa contraposée l'est aussi.}
\end{answers}
\begin{explanations}
L'implication réciproque est "$x>0\implies x^2>0$" (ce qui est une affirmation vraie). \\
La contraposée est "$x\le 0\implies x^2\le 0$" (ce qui est une affirmation fausse).\\
Une implication et sa contraposée sont des propositions équivalentes (donc vraies en même temps, et fausses en même temps). \\
\end{explanations}
\end{question}
}
